% Conformal prediction in failure attribution: case studies -- two-page reading report.
% Build: ./build.sh case_studies  (pdflatex + bibtex, output case_studies.pdf)
\documentclass[11pt,a4paper]{article}
\usepackage[margin=1.8cm,top=1.4cm,bottom=1.5cm,footskip=0.7cm]{geometry}
\usepackage[T1]{fontenc}
\usepackage[utf8]{inputenc}
\usepackage{amsmath,amssymb}
\usepackage{times}  % Times text; mathptmx needs the rsfs font, which this TeX install lacks
\usepackage{enumitem}
\usepackage{xcolor}
\usepackage[font=small,labelfont=bf,skip=4pt]{caption}
\usepackage[numbers,sort&compress]{natbib}
\usepackage[hidelinks]{hyperref}

\setlist{nosep,leftmargin=1.3em}
\setlength{\parskip}{2pt}
\setlength{\parindent}{0pt}
\setlength{\textfloatsep}{7pt plus 2pt}
\setlength{\abovedisplayskip}{4pt plus 1pt}
\setlength{\belowdisplayskip}{4pt plus 1pt}
\setlength{\abovedisplayshortskip}{2pt}
\setlength{\belowdisplayshortskip}{3pt}
\setlength{\emergencystretch}{2em}
\pagestyle{plain}
\makeatletter
\renewcommand\section{\@startsection{section}{1}{\z@}{6pt plus 2pt minus 1pt}{1.5pt}{\normalsize\bfseries}}
\let\oldbib\thebibliography
\renewcommand\thebibliography[1]{\oldbib{#1}\setlength{\itemsep}{0pt}\setlength{\parskip}{0pt}\footnotesize}
\makeatother

% ---------- notation ----------
\newcommand{\C}{\mathcal{C}}
\newcommand{\qhat}{\hat q}
\newcommand{\Prob}{\mathbb{P}}
\newcommand{\ystar}{y^{*}}
\newcommand{\jstar}{j^{*}}

% ---------- figure colours and helpers (plain picture environment) ----------
\definecolor{inset}{rgb}{0.23,0.30,0.70}
\definecolor{outset}{rgb}{0.80,0.82,0.86}
\definecolor{err}{rgb}{0.77,0.20,0.16}
\newcommand{\bbar}[4]{\put(#1,0){\textcolor{#4}{\rule{4mm}{#2mm}}}\put(#1,-3.6){\makebox(4,3){\scriptsize #3}}}
\newcommand{\cell}[4]{\put(#1,0){\textcolor{#3}{\rule{6.4mm}{5.4mm}}}\put(#1,0){\makebox(6.4,5.4){\scriptsize\textcolor{#4}{#2}}}}

\begin{document}

\begin{center}
{\large\bfseries Conformal Prediction in Failure Attribution: Case Studies}\\[2pt]
{\small Reading report \textperiodcentered\ November 2026}
\end{center}
\vspace{-5pt}

\textbf{Abstract.} Failure attribution methods for LLM multi-agent systems output one step, although the responsible step is often ambiguous and point predictions are usually wrong. Conformal prediction (CP) can instead return a set of steps with a statistical guarantee. We restate failure attribution as a set-prediction problem in CP notation, and then walk through the two works that apply CP to it, OAT and Conformal Agent Error Attribution. 
% For each we describe how steps are scored, how a threshold is calibrated on held-out runs, and what the resulting set is guaranteed to contain.

% =====================================================================
\section{Introduction}

\emph{Failure attribution} asks which step of a failed agent trajectory caused the failure \cite{zhang2025who}. Whether they prompt an LLM judge, run a structured reasoning pipeline or fine-tune a model, existing methods end by naming one step \cite{feng2026conformal}. Recent work questions this on two grounds. First, the label itself may not be a single step. MP-Bench \cite{mpbench2026} points out that benchmarks assume one deterministic failure step per run, while analysts who look at the same failure can reasonably disagree: one expected the agent to search the web first, another finds its choice to write code acceptable, and they blame different steps. Second, a single step says nothing about how far it can be trusted. Without a measure of uncertainty a practitioner cannot tell a reliable attribution from a guess \cite{feng2026conformal}, and guesses are common, as step-level accuracy on Who\&When is below 25\% \cite{zhang2025who}.

Conformal prediction (CP) is a natural answer to both concerns. It turns any step scorer into prediction sets that satisfy a guarantee at a level chosen by the user, in finite samples and without assumptions on the data distribution \cite{angelopoulos2021gentle}. This report reads the two works that apply CP to attribution \cite{yeh2026oat,feng2026conformal} with two questions in mind.

\textbf{Q1 (guarantee).} Does the method return a region of the trace that contains the decisive error with a user-specified probability?
\textbf{Q2 (usefulness).} How small is that region?

% =====================================================================
\section{Problem Setup}

We follow the sequential formulation of \cite{feng2026conformal}. A failed trajectory is a sequence of steps $x=(c_1,\dots,c_\ell)$ of variable length $\ell$. Its ground truth is a subset of steps $\ystar\subseteq x$: one decisive error in most datasets, several plausible steps in the view of \cite{mpbench2026}. Given calibration trajectories $(x_1,\ystar_1),\dots,(x_n,\ystar_n)$ and a miscoverage level $\alpha\in(0,1)$, the goal is a set predictor $\C(x)\subseteq x$ such that, for a new failed trajectory,
\begin{equation}
\Prob\big[\,\ystar_{n+1}\subseteq\C(x_{n+1})\,\big]\;\ge\;1-\alpha, \label{eq:cov}
\end{equation}
with $\C(x_{n+1})$ as small as possible. Split CP reaches \eqref{eq:cov} in three steps \cite{angelopoulos2021gentle}: (i) define a \emph{conformal score} $S(x,y)$, built from a heuristic scorer, that is large when $y$ is an implausible label for $x$; (ii) set the threshold $\qhat$ to the $\lceil(n+1)(1-\alpha)\rceil$-th smallest of the calibration scores $S(x_i,\ystar_i)$; (iii) output $\C(x)=\{y: S(x,y)\le\qhat\}$.

% =====================================================================
\section{Case Study 1: OAT \cite{yeh2026oat}}

\textbf{Formulation: outlier detection over steps.} OAT learns from successful runs only and has no labelled failures, so it cannot calibrate on the true error steps as in Section~2. It treats attribution as outlier detection instead: the steps of successful runs define what a \emph{normal} step looks like, and a step of a failed run is blamed when it does not look normal. CP has a variant for this setting \cite{angelopoulos2021gentle}. Given an outlier score and a sample of normal points, it picks a threshold that a new normal point exceeds with probability at most $\alpha$.

\textbf{Stage 1: anomaly scoring.} A frozen LLM encodes each step into a hidden state $h_t$. A neural controlled differential equation, trained only on successful runs, predicts the next state $\hat h_t$ from the preceding ones. The score is the prediction error $e_t=\lVert h_t-\hat h_t\rVert_2^2$: a step that departs from the usual flow of a successful run scores high.

\textbf{Stage 2: CP.} The calibration unit is a single step, not a run.
\begin{itemize}
\item \emph{Calibration scores.} Some successful runs are held out from training (20\% of the 103 runs on MCP-Atlas). The anomaly score $e_t$ is computed for every step of every held-out run, and all these scores are pooled into one list.
\item \emph{Threshold.} $\delta$ is the $(1-\alpha)$ quantile of the list, so a fraction $\alpha$ of the normal calibration steps score above it.
\item \emph{Prediction.} For a failed run, compute $e_t$ for each step and return the steps above the threshold, $\C_{\mathrm{OAT}}(x)=\{c_t: e_t>\delta\}$ (Figure~\ref{fig:cp}a).
\end{itemize}

\textbf{Guarantee.} A normal step is flagged with probability at most $\alpha$:
\begin{equation}
\Prob\big[\,e_t>\delta \;\big|\; c_t \text{ is a normal step}\,\big]\;\le\;\alpha. \label{eq:oat}
\end{equation}
In practice the flagged set $\C_{\mathrm{OAT}}(x)$ is reported as the failure-contributing steps of the run. Its size is not fixed: it can hold several steps or none. Since \eqref{eq:oat} limits false alarms and not misses, whether the set contains the error is measured empirically: at $\alpha=0.2$ it contains an annotated error in 57\% of MCP-Atlas runs and 33\% of Who\&When runs.

\begin{figure}[t]
\centering
\setlength{\unitlength}{1mm}
\begin{minipage}[b]{0.43\linewidth}\centering
\begin{picture}(66,29)(-3,-9)
\bbar{0}{3}{1}{outset}\bbar{6}{5}{2}{outset}\bbar{12}{4}{3}{outset}\bbar{18}{12}{4}{inset}\bbar{24}{6}{5}{outset}
\bbar{30}{3}{6}{outset}\bbar{36}{16}{\textcolor{err}{\textbf{7}}}{inset}\bbar{42}{6}{8}{outset}\bbar{48}{4}{\textcolor{err}{\textbf{9}}}{outset}\bbar{54}{5}{10}{outset}
\put(-1,0){\line(1,0){60}}
\multiput(-1,9)(2,0){30}{\line(1,0){1.1}}
\put(59.5,8.2){\scriptsize$\delta$}
\put(-3,16.5){\scriptsize$e_t$}
\put(1,13.5){\makebox(16,3)[r]{\scriptsize false alarm}}
\put(42,16.5){\makebox(14,3)[l]{\scriptsize hit}}
\put(40.5,9.8){\makebox(20,3){\scriptsize missed}}
\put(-1,-8){\makebox(60,3){\scriptsize step $t$ (errors in red); $\C_{\mathrm{OAT}}=\{c_4,c_7\}$}}
\end{picture}\\[1pt]
{\small (a) OAT: flag steps with $e_t>\delta$}
\end{minipage}\hfill
\begin{minipage}[b]{0.55\linewidth}\centering
\begin{picture}(84,29)(-14,-13)
\put(-14,8.2){\makebox(13,5.4)[l]{\scriptsize judge $p_i$}}
\put(-14,-0.3){\makebox(13,5.4)[l]{\scriptsize steps}}
\put(-14,-6.5){\makebox(13,4)[l]{\scriptsize $g(c_{j:\ell})$}}
% judge scores
\put(0,8.2){\makebox(6.4,5.4){\scriptsize .1}}\put(7,8.2){\makebox(6.4,5.4){\scriptsize .2}}\put(14,8.2){\makebox(6.4,5.4){\scriptsize .1}}\put(21,8.2){\makebox(6.4,5.4){\scriptsize .1}}\put(28,8.2){\makebox(6.4,5.4){\scriptsize\textbf{.7}}}
\put(35,8.2){\makebox(6.4,5.4){\scriptsize .2}}\put(42,8.2){\makebox(6.4,5.4){\scriptsize .1}}\put(49,8.2){\makebox(6.4,5.4){\scriptsize .1}}\put(56,8.2){\makebox(6.4,5.4){\scriptsize .2}}\put(63,8.2){\makebox(6.4,5.4){\scriptsize .1}}
% trajectory cells
\cell{0}{$c_1$}{outset}{black}\cell{7}{$c_2$}{outset}{black}\cell{14}{$c_3$}{outset}{black}\cell{21}{$c_4$}{inset}{white}\cell{28}{$c_5$}{err}{white}
\cell{35}{$c_6$}{inset}{white}\cell{42}{$c_7$}{inset}{white}\cell{49}{$c_8$}{inset}{white}\cell{56}{$c_9$}{inset}{white}\cell{63}{$c_{10}$}{inset}{white}
% suffix scores
\put(0,-6.5){\makebox(6.4,4){\scriptsize .19}}\put(7,-6.5){\makebox(6.4,4){\scriptsize .18}}\put(14,-6.5){\makebox(6.4,4){\scriptsize .16}}\put(21,-6.5){\makebox(6.4,4){\scriptsize\textbf{.15}}}\put(28,-6.5){\makebox(6.4,4){\scriptsize\textbf{.14}}}
\put(35,-6.5){\makebox(6.4,4){\scriptsize\textbf{.07}}}\put(42,-6.5){\makebox(6.4,4){\scriptsize\textbf{.05}}}\put(49,-6.5){\makebox(6.4,4){\scriptsize\textbf{.04}}}\put(56,-6.5){\makebox(6.4,4){\scriptsize\textbf{.03}}}\put(63,-6.5){\makebox(6.4,4){\scriptsize\textbf{.01}}}
\put(-14,-12){\makebox(84,3){\scriptsize $\qhat=0.15$: keep the longest suffix with $g\le\qhat$ (bold), so $\C_{\mathrm{LF}}=c_{4:10}$}}
\end{picture}\\[1pt]
{\small (b) Left filtration: error $c_5$ (red) is covered}
\end{minipage}
\caption{The calibrated threshold applied to a new failed run. (a) $\delta$ is a quantile of step scores from successful runs: it limits false alarms but can miss an error. (b) $\qhat$ is a quantile of one score per labelled failed run, $S_{\mathrm{LF}}=g(c_{\jstar:\ell})$ (here $0.14$).}\label{fig:cp}
\end{figure}

% =====================================================================
\section{Case Study 2: Conformal Agent Error Attribution \cite{feng2026conformal}}

\textbf{Formulation: contiguous sets.} The label is the single decisive error $\ystar=c_{\jstar}$, the earliest error the system cannot recover from. For debugging and for rolling the system back, a block of consecutive steps is preferred to a scattered set, so \eqref{eq:cov} becomes
\begin{equation}
\Prob\big[\,\ystar_{n+1}\in\C(x_{n+1})\,\big]\;\ge\;1-\alpha,\qquad \C(x_{n+1})=(c_j,\dots,c_k). \label{eq:contig}
\end{equation}
A vanilla baseline (split CP over steps as unordered classes) gives scattered sets; we do not detail it.

\textbf{Stage 1: scoring functions.} An LLM judge reads one step in its context and returns $p_i\in[0,1]$, its belief that the decisive error is in step $c_i$. Three judges are compared: a naive prompted gpt-4o-mini, an ensemble of four role-prompted judges whose scores are averaged, and a small model (Qwen3-1.7B) fine-tuned on trajectories with injected errors. Their step-level AUROC is 0.52, 0.55 and 0.76. A block of steps is scored by the normalised sum $g(c_{j:k})=\tfrac{1}{\ell}\sum_{i=j}^{k}p_i$, which never decreases when the block grows.

\textbf{Stage 2: CP by left (right) filtration.} Left filtration removes steps from the left until what remains scores below a threshold; it returns a suffix.
\begin{itemize}
\item \emph{Calibration score.} For a labelled failed run, the smallest threshold that keeps the error is the score of the suffix that starts at it: $S_{\mathrm{LF}}(x,\ystar)=g(c_{\jstar:\ell})$.
\item \emph{Threshold.} $\qhat$ is the $\lceil(n+1)(1-\alpha)\rceil$-th smallest of the $n$ calibration scores, one per run.
\item \emph{Prediction.} $\C_{\mathrm{LF}}(x)$ is the longest suffix $c_{j:\ell}$ with $g(c_{j:\ell})\le\qhat$ (Figure~\ref{fig:cp}b).
\end{itemize}
Right filtration mirrors this with prefixes, $S_{\mathrm{RF}}=g(c_{1:\jstar})$; two-way filtration intersects the two.

\textbf{Guarantee.} If calibration and test failures are exchangeable, then $1-\alpha\le\Prob[\ystar\in\C(x)]<1-\alpha+\tfrac{1}{n+1}$ for any judge. This is \eqref{eq:contig}, so Q1 is answered positively. On Q2: at $1-\alpha=0.8$ the best variant on Who\&When (right filtration) removes 31\% of the steps, so the set still holds about two thirds of the trajectory.

% =====================================================================
\section{Discussion}

\textbf{CP gives validity; the scorer gives usefulness.} The guarantee holds for any scorer, even a random one, because the threshold adapts until the guarantee is met. How small and how accurate the set is depends on the scorer alone \cite{angelopoulos2021gentle}, so attribution quality still rests on step scoring, and current scorers are weak. In \cite{feng2026conformal} the two prompted judges reach an AUROC of 0.52 and 0.55 (chance is 0.50) and a top-1 accuracy of 0.11 and 0.16 (random guessing is 0.12). OAT's flagged set recovers 48\% of the annotated error steps on MCP-Atlas and 33\% on Who\&When.

% \textbf{What the sets add on real failures is not yet shown.} On Who\&When, the only dataset of natural failures they use, \cite{feng2026conformal} report coverage and removal rate but no attribution accuracy; scorer accuracy is given on injected errors only. The best set there still keeps about two thirds of the trajectory, and it is not compared with point-prediction baselines. How much the method helps in practice therefore remains open.

\bibliographystyle{unsrtnat}
\bibliography{refs}
\end{document}
